\chapter*{Abstract}
\addcontentsline{toc}{chapter}{Abstract}

Automatic volcano-seismic monitoring requires detecting events in
continuous records and classifying them by type. Operational systems solve
both stages separately: a detector sets the event boundaries and a
classifier, trained on analyst-made segmentations, assigns the class. This
thesis develops an algorithm for volcano-seismic event identification that
adapts the classification stage to the variability of the window produced
by the detector, using records from the Nevados de Chillán Volcanic
Complex.

The work empirically evaluates three detectors (STA/LTA, PhaseNet and
EQTransformer) at two levels, event onset identification and quality of the
temporal delimitation; builds an adaptive windowing procedure fitted to the
actual duration of the events, whose range exceeds by two orders of
magnitude the fixed window commonly used; and validates the integrated
algorithm by contrasting alternative convolutional architectures and
quantifying the degradation between reference and automatic windows.

\vspace{1cm}
\noindent\textbf{Keywords:} volcano seismology, event detection, automatic
classification, convolutional neural networks, Nevados de Chillán.